% Draft text for thesis/Chapters/2.Introduction.tex, if the learned storage policy is adopted.
% Not in the thesis. Written 2026-10-22, see learned_storage_policy.md, "Definition 2026-10-22".
% [Axx] is a placeholder. The tag is numbered when the change is adopted and logged in GAPS.md.
%
% Part A applies either way. Parts B to D apply only if the learned policy is adopted.

% =================================================================
% Part A. Where the rules come from.
% sec:pf:storage, before "The two functional forms, the weights ..." (about line 442).
% =================================================================

% A1, the rules as they are now.
The rules follow the storage literature. Fixed storage is dedicated storage.
The demand rule is class-based storage by the cube-per-order index
\citep{heskett1963coi,dekoster2007picking}, the same kind of rule as the ABC
classes that \citet{WaubertdePuiseau2022DynamicSL} use as their baseline. The
congestion-aware rule adds the traffic estimates to the demand rule.

% A2, the same with the learned rule in place of the congestion-aware rule.
% The rules follow the storage literature. Fixed storage is dedicated storage.
% The demand rule is class-based storage by the cube-per-order index
% \citep{heskett1963coi,dekoster2007picking}, the same kind of rule as the ABC
% classes that \citet{WaubertdePuiseau2022DynamicSL} use as their baseline. The
% learned rule follows their method, a policy trained with deep reinforcement
% learning, and adds the traffic of the fleet to what it sees and to what it is
% judged by.

% =================================================================
% Part B. The learned rule.
% sec:pf:storage, after the paragraph that defines \mathcal{F} (about line 405).
% =================================================================

% [Axx]
A storage rule can also be learned. The learned rule places the arriving stock
one shelf at a time. It takes the SKUs in order of demand, as the demand rule
does. For each SKU it looks at the warehouse and chooses a shelf that still has
room, and as many units as fit go there. That shelf gives one put-away task. If
units are left, the rule chooses again among the shelves that still have room.
Units that fit on no shelf are not placed, as under every rule. If 30 units of
tea arrive and the rule chooses shelf $B$, which has room for 20, then 20 go to
$B$ and the rule chooses a shelf for the last 10.

The rule sees what the other rules read, the stock, the units on their way and
the three estimates, and the SKU it is placing. It chooses among the shelves
themselves, not among groups of shelves, so no division of the warehouse is
needed and the same rule runs on a warehouse of any size.

Formally, let $x^{\mathrm{in},(n)}$ be the units on their way after the first
$n$ choices of the epoch at $t$, with $x^{\mathrm{in},(0)} = x^{\mathrm{in},-}_t$,
and let $k_n$ be the SKU of choice $n$. The shelves with room for one more unit
of $k_n$ are
\[
  V^{\mathrm{room}}_n = \Bigl\{ v \in V_{\mathrm{str}} :
    b_{k_n} + \sum_{k \in \mathcal{K}} b_k
      \bigl( x_t^{-}(k,v) + x^{\mathrm{in},(n-1)}(k,v) \bigr) \le \mathrm{cap}(v) \Bigr\},
\]
and choice $n$ draws a shelf from the \emph{learned storage policy},
\begin{equation}
  v_n \sim \pi^{\mathrm{str}}_\theta\bigl(\,\cdot \mid x_t^{-},\, x^{\mathrm{in},(n-1)},\,
    k_n,\, \hat\rho_t,\, \hat\mu_t,\, \hat w_t \bigr),
  \qquad v_n \in V^{\mathrm{room}}_n .
  \label{eq:storagepolicy}
\end{equation}
As many units of $k_n$ as fit on $v_n$ are added to $x^{\mathrm{in},(n)}$. Every
choice stays within capacity, moves nothing and adds only arriving units, so the
result meets the three conditions of \nref{eq:storageupdate}. The draws come
from the run's seed, so for each seed the learned rule $F_{\mathrm{lrn}}$ is one
element of $\mathcal{F}$.

% =================================================================
% Part C. What the learned rule is trained for.
% sec:pf:game, after "The reward R_i itself is left open here ..." (about line 1050).
% =================================================================

% [Axx]
A learned storage rule has a learning problem of its own. It decides on the
slow clock, once per put-away, and the routing controller is part of the world
it learns in. While the storage rule learns, the controller is held fixed, so
the controller's choices belong to the transition kernel of the storage rule,
just as the storage update belongs to the kernel $P$ of the routers.

The storage rule is judged by the loaded journeys its choices settle. A shelf
choice decides two journeys, the put-away from the receiving side to the shelf
and each pick from the shelf to its delivery vertex. The time a robot spends on
them while carrying the units, moving or blocked, is the part of the service
time a shelf choice controls. The wait before assignment and the empty run to
the shelf depend on where the fleet happens to be, so they are left out. The
trip distance $d^{\mathrm{trip}}$~\eqref{eq:tripdistance} is the same journey in
an empty warehouse. The loaded time adds what traffic costs on that journey,
which the trip distance leaves out.

Formally, the \emph{loaded time} of task $\tau_j$ is
$\zeta^{\mathrm{load}}_j = t_j^{\mathrm{finish}} - t_j^{\mathrm{pickup}}$. With
$\mathcal{C}^{\mathrm{put}}_T$ the put-aways completed by $T$, the storage
policy is trained for
\begin{equation}
  \min_{\theta}\; \mathbb{E}_{\pi^{\mathrm{str}}_\theta} \Bigl[
    \sum_{\tau_j \in \mathcal{C}^{\mathrm{put}}_T} \zeta^{\mathrm{load}}_j
    + \sum_{\tau_j \in \mathcal{C}^{\mathrm{pick}}_T} \zeta^{\mathrm{load}}_j \Bigr].
  \label{eq:storagelearning}
\end{equation}
Both sums are needed. Orders are the same under every rule, so every pick
happens whatever the rule does, and only its loaded time depends on the shelf.
The number of put-aways does depend on the rule, because a SKU spread over more
shelves takes more trips, and the first sum charges for each. How each term is
credited to the choice that caused it is given in Section~\ref{sec:method:storage}.

% =================================================================
% Part D. Rows for the notation table (Appendices/A.Notation.tex).
% =================================================================
% F_fix, F_dem, F_cng row becomes F_fix, F_dem, F_lrn if the learned rule replaces F_cng.
%
% \( \pi^{\mathrm{str}}_\theta \) & \nameeq{eq:storagepolicy} & Learned storage policy, which draws a shelf with room for each put-away \\
% \( F_{\mathrm{lrn}} \) & learned storage rule & The storage rule that places arriving stock by \( \pi^{\mathrm{str}}_\theta \) \\
% \( \zeta^{\mathrm{load}}_j \) & loaded time & \( t_j^{\mathrm{finish}} - t_j^{\mathrm{pickup}} \), the time task \( \tau_j \)'s units are carried \\
% \( \mathcal{C}^{\mathrm{put}}_T \) & completed put-aways & Put-away tasks completed by \( T \) \\
%
% Helpers local to eq:storagepolicy, not in the table: n, k_n, v_n, x^{in,(n)}, V^room_n.
